% !TEX root = task-3-19-1.tex
% Editable task bank. Append new BankTask blocks at the end of the document.
\documentclass[12pt,a4paper]{article}
\usepackage{physics-textbook}
\usepackage{graphicx,needspace}
\usetikzlibrary{decorations.pathmorphing}
\newif\ifShowSolutions
\ifdefined\TutorEdition
  \ShowSolutionstrue
\else
  \ShowSolutionsfalse % Change to \ShowSolutionstrue to display solutions.
\fi
\newcounter{banktask}
\newsavebox{\BankTaskBox}
\newcommand{\BankTaskContent}[3]{%
  \noindent\textbf{\thebanktask.}\quad #2\par
  \ifShowSolutions
    \Needspace{50mm}% Keep the solution heading with its opening text or diagram.
    #3\par
    \ifdefined\TutorEdition
      \ifcsname TutorSource#1\endcsname
        \nopagebreak\smallskip{\small\csname TutorSource#1\endcsname\par}%
      \fi
    \fi
  \fi
}
\newcommand{\BankTask}[3]{%
  \refstepcounter{banktask}\par\medskip
  \sbox{\BankTaskBox}{\begin{minipage}{\linewidth}%
    \BankTaskContent{#1}{#2}{#3}\end{minipage}}%
  \ifdim\dimexpr\ht\BankTaskBox+\dp\BankTaskBox\relax>\dimexpr\textheight-20mm\relax
    \Needspace{35mm}\BankTaskContent{#1}{#2}{#3}%
  \else
    \Needspace{\dimexpr\ht\BankTaskBox+\dp\BankTaskBox\relax}%
    \noindent\usebox{\BankTaskBox}\par
  \fi
  \medskip
}
\begin{document}
\BookSubsection{3.19.1}{Магнитный поток}


% Source 9C35AF: https://ege.fipi.ru/bank/questions.php?search=1&proj=BA1F39653304A5B041B656915DC36B38&qid=9C35AF
\BankTask{9C35AF}{
В опыте по наблюдению электромагнитной индукции квадратная рамка из одного витка тонкого провода находится в однородном магнитном поле, перпендикулярном плоскости рамки. Индукция магнитного поля равномерно возрастает от $0$ до максимального значения $B_{\text{макс}}$ за время $T$. При этом в рамке возбуждается ЭДС индукции, равная $6\,\text{мВ}$. Какая ЭДС индукции (в $\text{мВ}$) возникнет в рамке, если $T$ уменьшить в $3$ раза, а $B_{\text{макс}}$ уменьшить в $2$ раза?
}{
\smallskip\noindent\textbf{Решение от GPT}\par
Модуль ЭДС равен $|\mathcal{E}|=S B_{\text{макс}}/T$. Поэтому
\[
|\mathcal{E}'|=\frac{S(B_{\text{макс}}/2)}{T/3}=\frac32|\mathcal{E}|=9\,\text{мВ}.
\]
Ответ: $9\,\text{мВ}$.
}

% Source Fe341D: https://ege.fipi.ru/bank/questions.php?search=1&proj=BA1F39653304A5B041B656915DC36B38&qid=Fe341D
\BankTask{Fe341D}{
Проволочная рамка площадью $50\,\text{см}^{2}$ помещена в однородное магнитное поле так, что плоскость рамки перпендикулярна вектору индукции $\vec B$. Проекция $B_n$ индукции магнитного поля на нормаль к плоскости рамки изменяется во времени $t$ согласно графику на рисунке.
\begin{enumerate}
\item Модуль ЭДС индукции, возникающей в рамке в промежутке времени от $5$ до $8\,\text{мс}$, равен $4\,\text{В}$.
\item Направления индукционного тока в рамке в промежутках времени от $0$ до $2\,\text{мс}$ и от $5$ до $8\,\text{мс}$ совпадают.
\item Модуль изменения магнитного потока в рамке минимален в промежутке времени от $2$ до $4\,\text{мс}$.
\item Модуль силы индукционного тока в рамке максимален в промежутке времени от $0$ до $2\,\text{мс}$.
\item ЭДС индукции в рамке отлична от $0$ только в промежутках времени от $0$ до $2\,\text{мс}$ и от $5$ до $8\,\text{мс}$.
\end{enumerate}
\begin{center}
\begin{tikzpicture}[book diagram,x=0.65cm,y=0.65cm]
\draw[step=1,gray,very thin] (0,-2) grid (8,3);
\draw[-{book tip}] (0,-2.25)--(0,4.2);
\draw[-{book tip}] (0,0)--(9.15,0);
\node[above left] at (0,3.85) {$B_n,\,\text{Тл}$};
\node[below right] at (8.85,-0.05) {$t,\,\text{мс}$};
\foreach \x in {1,2,3,4,5,6,7,8}{\node[below,fill=white,inner sep=1pt] at (\x,0) {$\x$};}
\foreach \y in {-2,-1,0,1,2,3}{\node[left] at (0,\y) {$\y$};}
\draw[thick] (0,-2)--(2,2)--(4,-1)--(5,-1)--(8,3);
\end{tikzpicture}
\end{center}
}{
\smallskip\noindent\textbf{Решение от GPT}\par
По закону Фарадея $|\mathcal E|=S|\Delta B_n/\Delta t|$, где $S=5\cdot10^{-3}\,\text{м}^2$; на участках $0$–$2$, $2$–$4$, $4$–$5$ и $5$–$8\,\text{мс}$ получаем соответственно $10$, $7{,}5$, $0$ и $20/3\,\text{В}$, поэтому утверждение 1 неверно. На первом и последнем участках $B_n$ возрастает, следовательно, направления индукционного тока совпадают: утверждение 2 верно. Минимальное изменение потока равно нулю на участке $4$–$5\,\text{мс}$, поэтому утверждение 3 неверно. При постоянном сопротивлении $|I|=|\mathcal E|/R$ максимально на участке $0$–$2\,\text{мс}$: утверждение 4 верно. На участке $2$–$4\,\text{мс}$ ЭДС также отлична от нуля, поэтому утверждение 5 неверно. Ответ: 2, 4.
}

\end{document}
